\begin{longtable}{|p{\textwidth}|}
\caption{Example CDL description of 1D dataset}
\label{tab:cdl-1D} \\
\hline \endhead
\hline \endfoot
netcdf 1D example\\
dimensions\\
\hline
\rowcolor{YellowGreen}  time = 9497\\
\hline

coordinates\\
\hline
\rowcolor{Apricot}\hspace{0.5cm}int32 time (time)\\
\rowcolor{Apricot}\hspace{0.5cm}\hspace{0.5cm}time:long\_name = "center time of averaging period"\\
\rowcolor{Apricot}\hspace{0.5cm}\hspace{0.5cm}time:axis = "T"\\
\rowcolor{Apricot}\hspace{0.5cm}\hspace{0.5cm}time:bounds = "time\_bnds"\\
\rowcolor{Apricot}\hspace{0.5cm}\hspace{0.5cm}time:comment = ""\\
\rowcolor{Apricot}\hspace{0.5cm}\hspace{0.5cm}time:coverage\_content\_type = "coordinate"\\
\rowcolor{Apricot}\hspace{0.5cm}\hspace{0.5cm}time:standard\_name = "time"\\
\rowcolor{Apricot}\hspace{0.5cm}\hspace{0.5cm}time:units = "hours since 1992-01-01T12:00:00"\\
\rowcolor{Apricot}\hspace{0.5cm}\hspace{0.5cm}time:calendar = "proleptic\_gregorian"\\
\hline

data variables\\
\hline
\hspace{0.5cm}float32 global\_mean\_barystatic\_sea\_level\_anomaly (time)\\
\hspace{0.5cm}\hspace{0.5cm}global\_mean\_barystatic\_sea\_level\_anomaly:\_FillValue = "9.969209968386869e+36"\\
\hspace{0.5cm}\hspace{0.5cm}global\_mean\_barystatic\_sea\_level\_anomaly:coverage\_content\_type = "modelResult"\\
\hspace{0.5cm}\hspace{0.5cm}global\_mean\_barystatic\_sea\_level\_anomaly:long\_name = "Global mean of barystatic sea level anomaly"\\
\hspace{0.5cm}\hspace{0.5cm}global\_mean\_barystatic\_sea\_level\_anomaly:standard\_name = ""\\
\hspace{0.5cm}\hspace{0.5cm}global\_mean\_barystatic\_sea\_level\_anomaly:units = "m"\\
\hspace{0.5cm}\hspace{0.5cm}global\_mean\_barystatic\_sea\_level\_anomaly:comment = "Global mean barystatic sea level anomaly due to changes in total ocean mass. Note: ECCOv4 uses a volume-conserving Boussinesq formulation of the MITgcm with a free-surface boundary condition with real freshwater flux forcing. Changes in ocean mass due to evaporation, precipitation, runoff, and sea-ice growth/melt are reflected in model sea level. However, as a consequence of the Boussinsq formulation, changes to seawater density due to net buoyancy fluxes (e.g., global mean surface heating/cooling) do not change model sea level anomaly (ETAN) via seawater expansion/contraction. Changes in global ocean density therefore induce a spurious change in model ocean bottom pressure (PHIBOT) via 'virtual mass fluxes'. The 'Greatbatch correction' is a time varying, globally-uniform correction to account for changes in global mean density in Boussinesq models. This correction is used to calculate dynamic sea surface height (SSH) and ocean bottom pressure (OBP). Importantly, there is no dynamical significance to the Greatbatch correction but it is required to account for steric changes in global sea level. See Greatbatch, 1994. J. of Geophys. Res. Oceans, doi.org/10.1029/94JC00847"\\
\hspace{0.5cm}\hspace{0.5cm}global\_mean\_barystatic\_sea\_level\_anomaly:valid\_min = "-0.045110903680324554"\\
\hspace{0.5cm}\hspace{0.5cm}global\_mean\_barystatic\_sea\_level\_anomaly:valid\_max = "0.0434933640062809"\\
\hspace{0.5cm}\hspace{0.5cm}global\_mean\_barystatic\_sea\_level\_anomaly:coordinates = "time"\\
\hspace{0.5cm}float32 global\_mean\_sea\_level\_anomaly (time)\\
\hspace{0.5cm}\hspace{0.5cm}global\_mean\_sea\_level\_anomaly:\_FillValue = "9.969209968386869e+36"\\
\hspace{0.5cm}\hspace{0.5cm}global\_mean\_sea\_level\_anomaly:coverage\_content\_type = "modelResult"\\
\hspace{0.5cm}\hspace{0.5cm}global\_mean\_sea\_level\_anomaly:long\_name = "Global mean of dynamic SSH"\\
\hspace{0.5cm}\hspace{0.5cm}global\_mean\_sea\_level\_anomaly:standard\_name = ""\\
\hspace{0.5cm}\hspace{0.5cm}global\_mean\_sea\_level\_anomaly:units = "m"\\
\hspace{0.5cm}\hspace{0.5cm}global\_mean\_sea\_level\_anomaly:comment = "Global mean of dynamic sea level anomaly, equivalent to global mean sea level change. Note: ECCOv4 uses a volume-conserving Boussinesq formulation of the MITgcm with a free-surface boundary condition with real freshwater flux forcing. Changes in ocean mass due to evaporation, precipitation, runoff, and sea-ice growth/melt are reflected in model sea level. However, as a consequence of the Boussinsq formulation, changes to seawater density due to net buoyancy fluxes (e.g., global mean surface heating/cooling) do not change model sea level anomaly (ETAN) via seawater expansion/contraction. Changes in global ocean density therefore induce a spurious change in model ocean bottom pressure (PHIBOT) via 'virtual mass fluxes'. The 'Greatbatch correction' is a time varying, globally-uniform correction to account for changes in global mean density in Boussinesq models. This correction is used to calculate dynamic sea surface height (SSH) and ocean bottom pressure (OBP). Importantly, there is no dynamical significance to the Greatbatch correction but it is required to account for steric changes in global sea level. See Greatbatch, 1994. J. of Geophys. Res. Oceans, doi.org/10.1029/94JC00847"\\
\hspace{0.5cm}\hspace{0.5cm}global\_mean\_sea\_level\_anomaly:valid\_min = "-0.055836163461208344"\\
\hspace{0.5cm}\hspace{0.5cm}global\_mean\_sea\_level\_anomaly:valid\_max = "0.0552055686712265"\\
\hspace{0.5cm}\hspace{0.5cm}global\_mean\_sea\_level\_anomaly:coordinates = "time"\\
\hspace{0.5cm}float64 global\_mean\_sterodynamic\_sea\_level\_anomaly (time)\\
\hspace{0.5cm}\hspace{0.5cm}global\_mean\_sterodynamic\_sea\_level\_anomaly:\_FillValue = "9.969209968386869e+36"\\
\hspace{0.5cm}\hspace{0.5cm}global\_mean\_sterodynamic\_sea\_level\_anomaly:coverage\_content\_type = "modelResult"\\
\hspace{0.5cm}\hspace{0.5cm}global\_mean\_sterodynamic\_sea\_level\_anomaly:long\_name = "Global mean of sterodynamic sea level anomaly"\\
\hspace{0.5cm}\hspace{0.5cm}global\_mean\_sterodynamic\_sea\_level\_anomaly:standard\_name = ""\\
\hspace{0.5cm}\hspace{0.5cm}global\_mean\_sterodynamic\_sea\_level\_anomaly:units = "m"\\
\hspace{0.5cm}\hspace{0.5cm}global\_mean\_sterodynamic\_sea\_level\_anomaly:comment = "Steric sea level anomaly associated with seawater expansion/contraction due to density changes. Note: ECCOv4 uses a volume-conserving Boussinesq formulation of the MITgcm with a free-surface boundary condition with real freshwater flux forcing. Changes in ocean mass due to evaporation, precipitation, runoff, and sea-ice growth/melt are reflected in model sea level. However, as a consequence of the Boussinsq formulation, changes to seawater density due to net buoyancy fluxes (e.g., global mean surface heating/cooling) do not change model sea level anomaly (ETAN) via seawater expansion/contraction. Changes in global ocean density therefore induce a spurious change in model ocean bottom pressure (PHIBOT) via 'virtual mass fluxes'. The 'Greatbatch correction' is a time varying, globally-uniform correction to account for changes in global mean density in Boussinesq models. This correction is used to calculate dynamic sea surface height (SSH) and ocean bottom pressure (OBP). Importantly, there is no dynamical significance to the Greatbatch correction but it is required to account for steric changes in global sea level. See Greatbatch, 1994. J. of Geophys. Res. Oceans, doi.org/10.1029/94JC00847"\\
\hspace{0.5cm}\hspace{0.5cm}global\_mean\_sterodynamic\_sea\_level\_anomaly:valid\_min = "-0.017658796143049296"\\
\hspace{0.5cm}\hspace{0.5cm}global\_mean\_sterodynamic\_sea\_level\_anomaly:valid\_max = "0.017642477223663407"\\
\hspace{0.5cm}\hspace{0.5cm}global\_mean\_sterodynamic\_sea\_level\_anomaly:coordinates = "time"\\
\hline
\end{longtable}